
\subsubsection{2.1 Equivariant Weight-Space Encoders}

The permutation symmetry of neural network weights motivates encoder architectures that respect this symmetry by construction. \textbf{DWSNets} \citep{Navon2023Equivariant} introduced the first permutation-equivariant layers for homogeneous MLP weight spaces; the authors report approximately $R^{2}\approx 0.89$ for accuracy prediction on a private MNIST model zoo. \textbf{GNN-NFN} \citep{Kofinas2024Graph} extended this approach to heterogeneous architectures by treating each neural network as a computational graph (nodes = neurons, edges = weight matrices) and applying graph neural network operations equivariant to each layer's permutation group. GNN-NFN achieves state-of-the-art property prediction across multiple architecture families. \textbf{NFN} \citep{Zhou2023Permutation} and its universal extension \citep{Zhou2024Universal} auto-construct equivariant layers for arbitrary weight spaces. \textbf{Monomial-NFN} \citep{Tran2024Monomial} extends symmetry handling to scaling and sign-flipping, achieving a 34\% parameter reduction via theory-guided construction.

None of these works conduct a controlled comparison with plain encoders on shared train/test splits, and none report sample efficiency curves at systematically varied training set sizes. Performance is measured at full-data scale on custom or private splits, preventing direct comparison with plain encoder baselines.

\subsubsection{2.2 Plain Weight-Space Learning and Augmentation}

The ModelZooDataset [Schürholt et al., 2022] provides the standardized model zoos used in the present study. \textbf{Schürholt et al. [2021]} demonstrated that self-supervised representation learning on flattened weights achieves $R^{2}\approx 0.83$ for accuracy prediction on the ModelZooDataset MNIST zoo, establishing plain encoders as competitive baselines at full data. Meynent et al. [2025] showed that combining behavioral and structural signals outperforms structure alone.

Permutation augmentation --- randomly permuting neuron orderings during training --- is analogous to random crops in image learning: it provides some but not all benefits of architectural invariance. Its role as an intermediate condition between plain and equivariant encoders has been discussed conceptually but not tested in a controlled weight-space learning study. These plain encoder works share the limitation that no equivariant encoder comparison on the same shared splits was conducted.

\subsubsection{2.3 Expressivity Theory and Sample Complexity}

\textbf{Dayan, Eitan, and Maron [2026]} proved that all permutation-equivariant weight-space networks are equivalent in expressivity. Their theorem implies equivariant and plain encoders share the same function class ceiling when data is abundant. Expressivity equivalence does not imply sample complexity equivalence: a symmetry-consistent hypothesis space may be traversed with fewer examples even if ultimately equally expressive, analogous to the well-documented advantage of convolutional networks over plain MLPs on images at finite sample sizes.

The present work provides the empirical counterpart: it measures the practical sample complexity gap and shows it closes at full data (Section 5.4). \textbf{Herrmann et al. [2024]} contributed RNN model zoo datasets, finding functionalist representations outperform mechanistic ones --- a parallel to the finding that data-level augmentation can outperform structural encoding at very small zoo sizes.

\subsubsection{2.4 Positioning}

The present work is the first to: (1) compare equivariant and plain weight-space encoders on shared ModelZooDataset splits; (2) measure sample efficiency via systematic training size ablation; (3) include permutation augmentation as an explicit intermediate condition; and (4) report a data-regime crossover that qualifies the assumption that equivariant inductive bias is uniformly beneficial at low data.

